%% toc: P15 172 172.11,172.12
\chapter{Robust Conclusions, Causal Effects and Robust Decisions}
\label{sec:p15-w02b-p15-070}

The preceding chapters supply three separate kinds of robustness result.
An error interval decides a sign or a threshold comparison only when the
interval misses the threshold
(\cref{thm:p15-approximation-certificates,thm:p15-approximation-rational}).
A variance ranking or a threshold status can reverse when an input range
changes (\cref{thm:p15-global-ranking}). Huber's clipped location estimate
has a bounded population displacement under asymmetric contamination
(\cref{thm:p15-huber-asymmetric-bias}) and a minimax limiting risk under
symmetric contamination (\cref{thm:p15-huber-location-minimax}).
This chapter states the single notion of which these are instances: a
conclusion is robust when it holds at every member of a declared family of
models, input ranges, contaminations and approximation errors.
It then separates three claims that are often all called robustness:
a property of a statistical procedure over a neighborhood of laws, the value
of a causal effect, and the persistence of a recommended action.

The inputs throughout are a declared family, a real quantity or several
real scores defined at each member, and a threshold or ordering of interest.
The outputs are a three-valued verdict (established throughout the family,
reversed throughout the family, or unresolved), the margin by which an
established verdict survives, and, for a decision, the actions that remain
optimal at every member.
\index{robustness!of a conclusion}\index{robust decision}
\index{sensitivity analysis!robust conclusion}

\section{Robustness of signs, rankings and threshold conclusions}
\label{sec:p15-w02b-p15-070-signs}

Saltelli and Annoni describe a global analysis as one that shows a desired
inference persisting while the input assumptions vary over plausible
ranges; following Leamer, they credit a conclusion only when the range of
assumptions is wide enough to be credible and the resulting range of
inferences is narrow enough to be useful
\cite[Sec.~2, p.~1509]{SaltelliAnnoni2010Sensitivity}.
The definitions below make both halves of that requirement explicit.
The family is an input that must be justified; the verdict is a
mathematical consequence of the family.

\begin{hbdefinition}{def:p15-w02b-p15-070-robust-conclusion}{Declared families, conclusions and robust status}
A \emph{declared family} is a nonempty set \(\Theta\). Each member
\(\theta\in\Theta\) fixes everything the conclusion depends on: a model
variant, an input law or range, a contaminating distribution, an
approximation error, or a combination of these. The family is part of
the claim; enlarging or shrinking it gives a different claim.

A \emph{threshold conclusion} is given by a function
\(q:\Theta\to\mathbb R\) and a threshold \(c\in\mathbb R\). Its
\emph{status} at \(\theta\) is \(q(\theta)>c\), \(q(\theta)<c\) or
\(q(\theta)=c\). A \emph{sign conclusion} is the case \(c=0\).
A \emph{ranking conclusion} is given by scores
\(q_1,\ldots,q_m:\Theta\to\mathbb R\), \(m\geq2\), and a permutation
\(\sigma\) of \(\{1,\ldots,m\}\); it asserts
\(q_{\sigma(1)}(\theta)>q_{\sigma(2)}(\theta)>\cdots>q_{\sigma(m)}(\theta)\).

For a function \(D:\Theta\to\mathbb R\), the strict sign conclusion
\(D>0\) is \emph{robust over \(\Theta\)} when \(D(\theta)>0\) for every
\(\theta\in\Theta\); similarly for \(D<0\). The \emph{verdict} of
\(D\) over \(\Theta\) is \emph{established} if \(D>0\) is robust,
\emph{reversed} if \(D<0\) is robust, and \emph{unresolved} otherwise.
A threshold conclusion \(q>c\) is robust when \(D=q-c\) is robustly
positive; a ranking is robust when it holds at every member.

When the families form a nested collection \(\Theta_r\), \(r\geq0\),
with \(\Theta_r\subseteq\Theta_s\) for \(r\leq s\), the
\emph{robustness radius} of a conclusion is the supremum of the \(r\)
for which it is robust over \(\Theta_r\) (zero if there is none).
By nesting, the conclusion is robust over every \(\Theta_r\) with \(r\)
below the radius and over none with \(r\) above it.
\end{hbdefinition}

A robust conclusion is not a probability statement about the members of
\(\Theta\). No weighting of members is used, and a conclusion robust over
\(\Theta\) asserts nothing about a member outside \(\Theta\).
When a probability law on the family is also justified, the share of
members at which a conclusion holds is a different quantity, of the kind
reported in \cref{thm:p15-global-ranking}.

\begin{hbtheorem}{thm:p15-w02b-p15-070-reduction}{proved=annex:p15-w02b-p15-070-reduction}{Signs, thresholds and rankings reduce to signs of differences}
Let \(\Theta\) be a declared family and \(D:\Theta\to\mathbb R\).
\begin{enumerate}
\item If \(\inf_\Theta D>0\), then \(D>0\) is robust. If \(D>0\) is
      robust, then \(\inf_\Theta D\geq0\), and equality can occur:
      for \(\Theta=(0,1]\) and \(D(\theta)=\theta\), the conclusion is
      robust although the infimum is zero. If \(\Theta\) is a nonempty
      compact topological space and \(D\) is lower semicontinuous, the
      infimum is attained, and \(D>0\) is robust exactly when
      \(\min_\Theta D>0\).
\item Exactly one of the three verdicts holds. The verdict is
      unresolved exactly when some member has \(D(\theta)\leq0\) and
      some member has \(D(\theta')\geq0\); in particular it is
      unresolved whenever \(D\) takes both strict signs or vanishes at
      some member.
\item A ranking with scores \(q_1,\ldots,q_m\) and permutation
      \(\sigma\) is robust exactly when each of the \(m-1\) adjacent
      differences \(q_{\sigma(j)}-q_{\sigma(j+1)}\) is robustly positive.
\item For a pair \(i\ne j\), separation of the two ranges,
      \(\inf_\Theta q_i>\sup_\Theta q_j\), implies robust positivity of
      \(q_i-q_j\) but is not required for it: with
      \(\Theta=[0,10]\), \(q_1(\theta)=\theta+1\) and
      \(q_2(\theta)=\theta\), the ranges overlap and the ranking is
      robust. If instead \(\Theta=\Theta'\times\Theta''\) and
      \(q_i\) depends only on the first coordinate and \(q_j\) only on
      the second, then
      \(\inf_\Theta(q_i-q_j)=\inf_{\Theta'}q_i-\sup_{\Theta''}q_j\).
\end{enumerate}
\end{hbtheorem}

Item 4 records the computational trap in comparing two separately reported
intervals. When the sources of variation are shared, as with a common input
or a contaminant that affects both compared quantities, the difference must
be computed member by member. When the sources vary separately, the worst
case of the difference combines the two worst cases.

\begin{hbtheorem}{thm:p15-w02b-p15-070-margins}{proved=annex:p15-w02b-p15-070-margins}{Additive error margins and when they are sharp}
Let \(D:\Theta\to\mathbb R\) satisfy
\[
 D(\theta)=d_0+\sum_{r=1}^R e_r(\theta),\qquad
 |e_r(\theta)|\leq E_r\quad(\theta\in\Theta,\ 1\leq r\leq R),
\]
for a nominal value \(d_0\) and finite bounds \(E_r\geq0\).
Put \(E=\sum_rE_r\).
\begin{enumerate}
\item If \(d_0>E\), then \(D>0\) is robust; if \(d_0<-E\), then
      \(D<0\) is robust.
\item If some member has \(e_r(\theta)=-E_r\) for every \(r\)
      simultaneously, then \(d_0\leq E\) implies that \(D>0\) is not
      robust. This holds in particular when
      \(\Theta=\Theta_1\times\cdots\times\Theta_R\), each \(e_r\)
      depends only on the coordinate in \(\Theta_r\), and each attains
      \(-E_r\).
\item Without joint attainment the additive margin can be conservative:
      with \(\Theta=\{\theta_1,\theta_2\}\), \(R=2\),
      \(E_2=E_1>0\), \(e_1=(E_1,-E_1)\) and \(e_2=(-E_1,E_1)\), one has
      \(D=d_0\) at both members, so \(D>0\) is robust for every
      \(d_0>0\), including \(0<d_0\leq2E_1\).
\end{enumerate}
\end{hbtheorem}

\Cref{thm:p15-approximation-certificates} is the case of two approximated
terms. The theorem also covers a model difference evaluated with numerical
error, an input range and a contamination displacement, provided each
contribution has a bound valid throughout the declared family.
A margin computed from a bound that holds only at one member, or only in
a limit, does not enter item 1.

\begin{hbtheorem}{thm:p15-w02b-p15-070-linear-ranking}{proved=annex:p15-w02b-p15-070-linear-ranking}{Exact robust ranking of variance shares in a linear model}
Let \(Y=a_0+\sum_{i=1}^d a_iX_i\) with real coefficients and independent
square-integrable inputs of variances \(v_i\). Suppose the declared family
lets \((v_1,\ldots,v_d)\) range over a box
\(\prod_i[\ell_i,u_i]\), \(0\leq\ell_i\leq u_i<\infty\), on which
\(V=\sum_ia_i^2v_i>0\) throughout. Then, at each member,
\(S_i=S_{T_i}=a_i^2v_i/V\). For \(i\ne j\),
\[
\begin{aligned}
 S_i>S_j\text{ at every member}&\iff a_i^2\ell_i>a_j^2u_j,\\
 S_i\geq S_j\text{ at every member}&\iff a_i^2\ell_i\geq a_j^2u_j.
\end{aligned}
\]
For a level \(s_0\), \(S_i>s_0\) holds throughout the box exactly when
\[
 \frac{a_i^2\ell_i}{a_i^2\ell_i+\sum_{j\ne i}a_j^2u_j}>s_0,
\]
whose denominator is positive because it is \(V\) at a point of the box.
With \(a=(1,2)\) the ranking criterion is the condition
\(\ell_1\geq4u_2\) of \cref{thm:p15-global-ranking}.
\end{hbtheorem}

The model coefficients \(a_i\) are fixed in this statement. The
ranking of shares can reverse without any change in the coefficients,
because the shares measure contributions to output variance under the
declared input law, not derivatives of the model.

\subsection{Contamination displacement as the margin of a location sign}
\label{sec:p15-w02b-p15-070-contamination-margin}

Huber's starting point is that, under the indeterminacy model
\(F=(1-\varepsilon)\Phi+\varepsilon H\), the location parameter is not
uniquely determined unless further structure, such as symmetry, is added
\cite[Sec.~1, p.~73]{Huber1964RobustLocation}. His discussion of bias
treats asymmetric contamination of the normal component and gives
approximate upper bounds for the resulting displacement
\cite[Sec.~7, pp.~82--83]{Huber1964RobustLocation}. The next theorem
converts the exact displacement of
\cref{thm:p15-huber-asymmetric-bias} into the margin of a sign,
threshold or ranking conclusion about the center of the normal component.

\begin{hbdefinition}{def:p15-w02b-p15-070-asymmetric-family}{The asymmetric contamination family and its consistent centers}
Fix \(k>0\) and \(0\leq\varepsilon<1\). Let
\[
 \mathcal F_\varepsilon
 =\{(1-\varepsilon)\Phi+\varepsilon H:H\text{ a probability on }\mathbb R\}.
\]
For \(F\in\mathcal F_\varepsilon\), let \(c(F)\) be the unique zero of
the population score of \cref{thm:p15-huber-asymmetric-bias}, and let
\(c_*=c_*(\varepsilon)\) be the maximal displacement defined there for
\(\varepsilon<1/2\). Observations are \(X_i=\xi+W_i\) with
\(W_1,W_2,\ldots\) independent with law \(F\in\mathcal F_\varepsilon\),
and the clipped estimate \(T_{n,k}\) of
\cref{def:p15-huber-location-score} is computed with the known unit scale.
The \emph{population limit} is \(t(\xi,F)=\xi+c(F)\). For a real
\(t\), the \emph{set of consistent centers} is
\[
 \Xi_\varepsilon(t)=\{\xi\in\mathbb R:t(\xi,F)=t
                      \text{ for some }F\in\mathcal F_\varepsilon\}.
\]
\end{hbdefinition}

\begin{hbtheorem}{thm:p15-w02b-p15-070-contamination-margin}{proved=annex:p15-w02b-p15-070-contamination-margin}{The maximal displacement is the exact margin of a location conclusion}
In the setting of \cref{def:p15-w02b-p15-070-asymmetric-family}:
\begin{enumerate}
\item \(T_{n,k}\to t(\xi,F)\) almost surely.
\item For \(0\leq\varepsilon<1/2\) and every real \(t\),
      \(\Xi_\varepsilon(t)=[t-c_*,t+c_*]\). Every point of this interval
      is the center for some contaminant supported on two points.
\item Hence, for \(0\leq\varepsilon<1/2\) and a threshold \(\xi_0\), the conclusion \(\xi>\xi_0\)
      holds for every member consistent with the limit \(t\) exactly when
      \(t-\xi_0>c_*\); the conclusion \(\xi<\xi_0\) does so exactly
      when \(t-\xi_0<-c_*\); otherwise the verdict is unresolved.
      For two centers \(\xi_A,\xi_B\) observed through separately
      contaminated samples with limits \(t_A,t_B\), the ranking
      \(\xi_A>\xi_B\) is robust exactly when \(t_A-t_B>2c_*\).
      If both samples share one law \(F\), then
      \(t_A-t_B=\xi_A-\xi_B\) exactly.
\item The function \(\varepsilon\mapsto c_*(\varepsilon)\) on
      \([0,1/2)\) is continuous and strictly increasing, with
      \(c_*(0)=0\) and \(c_*(\varepsilon)\to\infty\) as
      \(\varepsilon\uparrow1/2\). For a margin \(\delta=t-\xi_0>0\),
      the conclusion \(\xi>\xi_0\) is robust over
      \(\mathcal F_\varepsilon\) exactly when
      \[
       \varepsilon<\varepsilon^*(\delta)\defeq
       \frac{-m_k(\delta)}{k-m_k(\delta)},
      \]
      and \(0<\varepsilon^*(\delta)<1/2\). This is its robustness
      radius in the nested families \(\mathcal F_\varepsilon\).
\item For \(1/2\leq\varepsilon<1\), every \(t\) and every \(\xi_0\),
      neither \(\xi>\xi_0\) nor \(\xi<\xi_0\) is robust.
\end{enumerate}
\end{hbtheorem}

The interval in item 2 is a statement about what the exact limit of this
estimator can reveal; it does not depend on the sample size. Sampling
error adds a second margin. If an interval \(J_n\) is shown to contain
\(t(\xi,F)\) with probability at least \(1-\alpha\) under every member of
the family, then on that event \(\xi\) lies in
\([\inf J_n-c_*,\sup J_n+c_*]\), so the same probability bound holds for
the enlarged interval. That uniform coverage is a separate hypothesis; the
pointwise normal limit of \cref{thm:p15-huber-normal-variance} does not
supply it, and \cref{thm:p15-huber-location-failure-cases} shows that
finite-sample suprema over contamination need not be controlled.

\begin{hbtheorem}{thm:p15-w02b-p15-070-linearized-margin}{proved=annex:p15-w02b-p15-070-linearized-margin}{Linearized displacement understates the margin}
For \(k>0\) and \(0<\varepsilon<1/2\),
\[
 c_*(\varepsilon)>\frac{\varepsilon k}{(1-\varepsilon)A(k)},
 \qquad A(k)=2\Phi(k)-1.
\]
Consequently, if a margin \(\delta=t-\xi_0\) satisfies
\(\varepsilon k/((1-\varepsilon)A(k))<\delta<c_*(\varepsilon)\), the
linearized value certifies \(\xi>\xi_0\) although some member of
\(\mathcal F_\varepsilon\) consistent with \(t\) has \(\xi=\xi_0\).
For \(\varepsilon=0.1\) and \(k=1.14\), the exact displacement satisfies
\(0.1703<c_*<0.1704\), whereas the linearized value is \(0.16986\)
to five decimals; the margin \(\delta=0.1702\) is certified by both the
linearized value and the rounded bound \(0.17\), but is not robust.
\end{hbtheorem}

Huber's display on p.~83 lists \(|c|\leq0.17\) for \(\varepsilon=0.1\),
\(k=1.14\), and calls the listed values approximate upper bounds obtained
by replacing a derivative at an intermediate point by its value at zero
\cite[Sec.~7, p.~83]{Huber1964RobustLocation}. The theorem shows the
direction of that approximation for the clipped normal score: it always
lowers the bound. For the two smaller contamination levels in the same
display, \(\varepsilon=0.01\), \(k=1.95\) and \(\varepsilon=0.001\),
\(k=2.64\), the exact values lie in \((0.020755,0.020765)\) and
\((0.0026645,0.0026650)\), below the listed \(0.021\) and \(0.0027\).
Rounding up kept those two entries above the exact values; the
linearization alone did not. A certification needs
an upper bound for the margin; an approximation of unknown direction
does not supply one (\cref{def:p15-approximation-relations}).

\begin{hbexample}{def:p15-w02b-p15-070-two-measurements}{Sign, threshold and ranking for two contaminated measurements}
Two quantities \(\xi_A,\xi_B\) are measured in units of the known normal
scale, each by its own sample, and each sample's error law may lie anywhere
in \(\mathcal F_\varepsilon\), separately for the two samples. Use
\(k=1\). Suppose the clipped estimates have limits \(t_A=0.35\) and
\(t_B=0.10\). Three conclusions are of interest: \(\xi_B>0\) (margin
\(0.10\)), \(\xi_A>0.2\) (margin \(0.15\)), and \(\xi_A>\xi_B\)
(difference \(0.25\), compared with \(2c_*\)).

By \cref{thm:p15-w02b-p15-070-contamination-margin}, the radii are
\(\varepsilon^*(0.10)\approx0.0638\) for the sign,
\(\varepsilon^*(0.15)\approx0.0927\) for the threshold, and, for the
ranking, the \(\varepsilon\) at which \(c_*=0.125\), namely
\(\varepsilon^*(0.125)\approx0.0785\). Evaluations of \(m_1\) give
\(0.0771<c_*(0.05)<0.0772\) and \(0.1632<c_*(0.1)<0.1633\).
\begin{center}
\begin{tabular}{lccc}
 & \(\xi_B>0\) & \(\xi_A>0.2\) & \(\xi_A>\xi_B\)\\\hline
 \(\varepsilon=0.05\) & robust & robust & robust\\
 \(\varepsilon=0.10\) & unresolved & unresolved & unresolved\\
 radius \(\varepsilon^*\) & \(0.0638\) & \(0.0927\) & \(0.0785\)
\end{tabular}
\end{center}
The three conclusions come from the same two limits but have different
radii. The ranking has the largest raw difference, yet its radius lies
between the other two, because separate contaminations of the two samples
can move the two limits in opposite directions, doubling the margin it
must exceed (\cref{thm:p15-w02b-p15-070-reduction}, item 4). If the
design guarantees that both samples share one error law, the ranking holds
at every \(\varepsilon<1\), since the common displacement cancels.
Reporting only that the estimator is \enquote{robust} would not indicate
which of these conclusions survives \(\varepsilon=0.1\): none does.
\end{hbexample}

\begin{hbwarning}{def:p15-w02b-p15-070-false-robustness}{Three ways to overstate robustness}
\emph{Checking selected members.} A conclusion verified at finitely
many members, or along each coordinate separately, need not hold
throughout the family. With \(\Theta=[-1,1]^2\) and
\(D(\theta)=1-2\theta_1\theta_2\), \(D=1\) on both coordinate axes,
yet \(D(1,1)=-1\). This is the interaction obstruction of
\cref{thm:p15-global-oat-obstruction} applied to a conclusion instead of
an output. Without a bound on \(D\) between the checked members, the
infimum in \cref{thm:p15-w02b-p15-070-reduction} has not been computed.

\emph{Leaving the family.} A verdict established over \(\Theta\) is
silent outside it. The ranking and threshold of
\cref{thm:p15-global-ranking} are robust while the range of \(X_1\)
stays \([-1,1]\) and reverse when that range becomes \([-3,3]\); the
location sign of the example above is robust at \(\varepsilon=0.05\)
and unresolved at \(\varepsilon=0.10\). A first-order approximation can
likewise decide a threshold near the expansion point and reverse it
farther out (\cref{thm:p15-approximation-rational}).

\emph{Using a margin that is not a bound.} A linearized, asymptotic or
rounded value of the worst-case displacement can lie below the true worst
case, as \cref{thm:p15-w02b-p15-070-linearized-margin} shows.
\end{hbwarning}

\begin{hbmethod}{def:p15-w02b-p15-070-certify}{Certifying a sign, ranking or threshold conclusion}
\begin{enumerate}
\item State the conclusion as a sign of a function \(D\) on the family:
      \(q-c\) for a threshold, adjacent differences for a ranking.
\item Declare the family: model variants, input ranges and laws,
      contamination neighborhoods, numerical and approximation errors.
      Record which sources of variation are shared by compared quantities
      and which vary separately.
\item Bound \(D\) below (and, for a reversed verdict, above) over the
      whole family, by an exact infimum, an attained minimum on a compact
      family, or additive margins valid at every member
      (\cref{thm:p15-w02b-p15-070-margins}). Use upper bounds for every
      margin; do not substitute a linearized or rounded value.
\item Report the verdict, the margin, and, for nested families, the
      radius. If the verdict is unresolved, report the members that
      realize both statuses and retain the comparison as unresolved
      rather than choosing the nominal answer.
\item For sampled estimates, add a sampling margin only when its coverage
      holds uniformly over the family, and state that hypothesis.
\end{enumerate}
\end{hbmethod}

\section{Distinguishing statistical robustness, causal effects and robust decisions}
\label{sec:p15-w02b-p15-070-three-claims}

Pearl draws the line between associational and causal concepts by asking
whether a concept can be defined from the joint distribution of observed
variables; causal concepts cannot, and claims involving them need causal
premises that observational data do not test
\cite[Secs.~2.1--2.2, pp.~99--100]{Pearl2009CausalOverview}.
He adds that sensitivity to causal assumptions, unlike sensitivity to a
statistical prior, does not shrink as the sample grows
\cite[Sec.~2.4, p.~101]{Pearl2009CausalOverview}.
Statistical robustness, as in Huber's analysis, belongs on the
associational side of that line: it concerns how a procedure behaves when
the law of the observations varies over a neighborhood. A decision adds a
third object, a loss that ranks actions. The three claims below have
different subjects and different failure modes.

\begin{hbdefinition}{def:p15-w02b-p15-070-three-claims}{Three different claims}
\emph{A statistical robustness claim} names a procedure, a target
functional of the observational law, a neighborhood of laws and a
performance criterion, and asserts a bound on the criterion over the
neighborhood. Examples are the exact maximal displacement \(c_*\) of
\cref{thm:p15-huber-asymmetric-bias} for the criterion
\(|c(F)|\) over \(\mathcal F_\varepsilon\), and the limiting risk
\(\mathcal R_\varepsilon\) of
\cref{def:p15-huber-symmetric-location-model}, minimized in
\cref{thm:p15-huber-location-minimax}. For a single law \(F\), write
\[
 \mathcal R(T;F)=\lim_{C\to\infty}\liminf_{n\to\infty}
                  E_F\min(nT_n^2,C^2),
\]
so that \(\mathcal R_\varepsilon(T)=\sup_{F\in\mathcal C_\varepsilon}
\mathcal R(T;F)\).

\emph{A causal-effect claim} names a model class and a quantity defined
from the intervention laws of each model
(\cref{def:p16-causal-intervention}), and asserts its value. The claim is
supported by observational data only to the extent that the quantity is
identified in the class (\cref{def:p16-causal-identification}).

\emph{A decision claim} names a nonempty set \(\mathcal A\) of actions,
a declared family \(\Theta\) and a loss
\(L:\mathcal A\times\Theta\to\mathbb R\) such that
\(L_*(\theta)=\min_{b\in\mathcal A}L(b,\theta)\) exists for every
\(\theta\). An action \(a\) is \emph{optimal at} \(\theta\) if
\(L(a,\theta)=L_*(\theta)\), and \emph{robustly optimal} if it is
optimal at every member. Its \emph{regret} is
\(r(a,\theta)=L(a,\theta)-L_*(\theta)\geq0\) and its \emph{maximal
regret} is \(\bar r(a)=\sup_\theta r(a,\theta)\). It is \emph{minimax}
if it minimizes \(\sup_\theta L(a,\theta)\in(-\infty,\infty]\) over
\(\mathcal A\), and \emph{minimax-regret} if it minimizes \(\bar r\).
A \emph{robust decision} is a robustly optimal action; when a decision
reports only a threshold status, it is robust when that status is robust
in the sense of \cref{def:p15-w02b-p15-070-robust-conclusion}.
\end{hbdefinition}

For expected-utility comparisons under a finite set of priors,
\cref{thm:p17-robust-dominance} already shows that comparison under every
prior is incomplete and that maxmin preference does not imply it. The
next theorem states the corresponding facts for a general loss, adds the
regret characterization and the dependence of minimax on how losses are
normalized, and connects a two-action decision with a threshold
conclusion.

\begin{hbtheorem}{thm:p15-w02b-p15-070-decisions}{proved=annex:p15-w02b-p15-070-decisions}{Robust decisions, regret and minimax risk}
In a decision claim as in \cref{def:p15-w02b-p15-070-three-claims}:
\begin{enumerate}
\item An action is robustly optimal exactly when its maximal regret is
      zero. If a robustly optimal action exists, the minimax-regret
      actions are exactly the robustly optimal actions.
\item Every robustly optimal action is minimax. The converse fails in two
      ways. With \(\Theta=\{\theta_1,\theta_2\}\) and loss vectors
      \(L(a,\cdot)=(0,5)\), \(L(b,\cdot)=(3,5)\), the action \(a\) is
      robustly optimal while \(b\) is also minimax. With
      \((0,10)\), \((10,0)\) and \((6,6)\), no action is robustly
      optimal and the unique minimax action is optimal at no member.
\item For \(\lambda>0\) and any \(h:\Theta\to\mathbb R\), the loss
      \(\lambda L(a,\theta)+h(\theta)\) has the same robustly optimal
      actions, regret \(\lambda r\), and the same minimax-regret actions.
      Its minimax actions can differ: with
      \(L(\mathrm{act},\cdot)=(-2,1)\), \(L(\mathrm{abstain},\cdot)=(0,0)\)
      the unique minimax action is to abstain, while after adding
      \(h=(2,-1)\) it is to act.
\item Let \(\mathcal A=\{\mathrm{act},\mathrm{abstain}\}\), and suppose
      that at each \(\theta\), acting is optimal exactly when
      \(q(\theta)\geq c\) and abstaining exactly when \(q(\theta)\leq c\),
      as for \(L(\mathrm{act},\theta)=c-q(\theta)\),
      \(L(\mathrm{abstain},\theta)=0\). Then acting is robustly optimal
      exactly when \(q\geq c\) on \(\Theta\), abstaining exactly when
      \(q\leq c\) on \(\Theta\), and a robust decision exists exactly when
      \(q-c\) does not take both strict signs. Acting is the unique
      optimal action at every member exactly when the threshold
      conclusion \(q>c\) is robust.
\end{enumerate}
\end{hbtheorem}

A minimax action is therefore a way of choosing when no robust decision
exists, and its choice depends on the loss normalization; a robust
decision, when it exists, is unaffected by that normalization.
Huber's estimator illustrates the first point within statistics itself.

\begin{hbtheorem}{thm:p15-w02b-p15-070-minimax-not-robust}{proved=annex:p15-w02b-p15-070-minimax-not-robust}{A minimax estimator need not be optimal at each member}
Let \(0<\varepsilon<1\) and \(k>0\). The normal law \(\Phi\) belongs to
\(\mathcal C_\varepsilon\). The sample mean \(\overline X_n\) and the
clipped estimate satisfy
\[
 \mathcal R(\overline X;\Phi)=1<
 \frac{B(k)}{A(k)^2}=\mathcal R(T_{\cdot,k};\Phi),
\]
with \(A,B\) as in \cref{thm:p15-huber-normal-variance}; for \(k=1\)
the right side is \(1.1073\) to four decimals. Meanwhile
\(\mathcal R_\varepsilon(\overline X)=\infty\) and, for the normalized
\(k\) of \cref{thm:p15-huber-least-favorable-normal},
\(\mathcal R_\varepsilon(T_{\cdot,k})=V^*(k,\varepsilon)<\infty\).
Thus the minimax estimator is not optimal at the member \(\Phi\), and the
estimator better at \(\Phi\) has unbounded worst-case risk.
\end{hbtheorem}

Minimax over a neighborhood is a statistical robustness claim about worst
cases. It is not the persistence claim of a robust decision, and neither
is a claim about a causal effect. The next theorem shows why no amount of
statistical robustness substitutes for identification.

\begin{hbtheorem}{thm:p15-w02b-p15-070-consistency-identification}{proved=annex:p15-w02b-p15-070-consistency-identification}{Consistency throughout a class requires identification}
Let \(\mathcal K\) be a class of models, each \(M\in\mathcal K\) with an
observational law \(P_M\) on a common measurable space
\((\mathcal O,\mathcal G)\), and let \(Q:\mathcal K\to\mathbb R\).
Call \(Q\) identified in \(\mathcal K\) if \(P_{M_1}=P_{M_2}\) implies
\(Q(M_1)=Q(M_2)\), as in \cref{def:p16-causal-identification}.
A procedure is a sequence of measurable maps
\(T_n:\mathcal O^n\to\mathbb R\) applied to independent observations
with law \(P_M\).
\begin{enumerate}
\item If \(P_{M_1}=P_{M_2}\), then for each \(n\) the law of \(T_n\) is
      the same under \(M_1\) and \(M_2\). Every property of the procedure
      defined from these laws and from a neighborhood of laws specified
      in terms of \(P_M\) (consistency for a constant, a limiting risk, a
      maximal displacement over the neighborhood) has the same value at
      \(M_1\) and \(M_2\).
\item If \(T_n\to Q(M)\) in probability under every \(M\in\mathcal K\),
      then \(Q\) is identified in \(\mathcal K\).
\item Consequently, in a class where \(Q\) is not identified, such as
      either class of \cref{thm:p16-causal-nonidentification}, no
      procedure is consistent for \(Q\) throughout the class, however
      robust it is.
\end{enumerate}
\end{hbtheorem}

\subsection{A worked case in which the three claims come apart}
\label{sec:p15-w02b-p15-070-worked-case}

\begin{hbdefinition}{def:p15-w02b-p15-070-confounded-shift}{A confounded shift model}
Let \(Z\) be a fair binary variable and \(W\) a real variable with law
\(F\), independent of \(Z\). For real \(\beta,\gamma\), the model
\(M(\beta,\gamma,F)\) has the structural assignments
\[
 X\defeq Z,\qquad Y\defeq\beta X+\gamma Z+W,
\]
with \(Z\) unobserved and \((X,Y)\) observed. Its intervention
\(\operatorname{do}(X=x)\) replaces only the assignment of \(X\) by the
constant \(x\in\{0,1\}\), keeping \(Z\), \(W\) and the assignment of
\(Y\), as in \cref{def:p16-causal-intervention}; write \(Y^{X=x}\) for
the resulting response. The \emph{shift effect} of \(X\) is a real
\(s\) with \(Y^{X=1}\) distributed as \(Y^{X=0}+s\). Put
\(\alpha=\beta+\gamma\).

The \emph{recorded arm laws} \(G_0,G_1\) are the laws of \(Y\) given
\(X=0\) and \(X=1\). A \emph{statistical contamination} of the records
replaces them by \(G_x=F_x(\cdot-\alpha x)\) with
\(F_0,F_1\in\mathcal F_\varepsilon\), for instance when recording errors
differ between the arms. It changes the recorded laws, not the structural
assignments. Each arm is estimated by the clipped estimate with \(k>0\),
applied to independent observations from \(G_x\), and the
\emph{clipped contrast} is the difference of the two arm estimates.
\end{hbdefinition}

\begin{hbtheorem}{thm:p15-w02b-p15-070-separation}{proved=annex:p15-w02b-p15-070-separation}{Robust estimation, causal effect and robust decision come apart}
In the confounded shift model:
\begin{enumerate}
\item The observational law of \((X,Y)\) gives \(X\) the fair
      binary law and, given \(X=x\), gives \(Y\) the law of
      \(\alpha x+W\). It depends on \((\beta,\gamma)\) only through
      \(\alpha\). The shift effect exists and equals \(\beta\). Hence
      the shift effect is not identified in any class containing
      \(M(\beta,\gamma,F)\) and \(M(\beta+\gamma,0,F)\) with
      \(\gamma\ne0\).
\item \emph{Robust, but for the observational contrast.} If
      \(F\in\mathcal F_\varepsilon\), \(0\leq\varepsilon<1\), and both
      arms are uncontaminated records of the model, the clipped contrast
      converges almost surely to \(\alpha\) for every contaminant \(H\).
      Under separate record contaminations with \(\varepsilon<1/2\), it
      converges to a limit \(\tau\in[\alpha-2c_*,\alpha+2c_*]\), and
      every point of this interval is attained. In the first case the
      limit equals the shift effect \(\beta\) exactly when \(\gamma=0\);
      in the second, the limit is \(\alpha\) plus a displacement fixed by
      the contaminants, and its distance from \(\beta\) depends on
      \(\gamma\), which the observational law does not reveal.
\item \emph{Identified, but not robustly estimated by every identifying
      functional.} In the subclass \(\gamma=0\), \(F=\Phi\), the shift
      effect equals \(\alpha\) and is identified by the difference of
      arm means and by the clipped contrast alike. Under contamination of
      the arm \(X=1\) alone, with \(0<\varepsilon<1/2\), the clipped
      contrast converges to a limit within \(c_*\) of \(\beta\), whereas
      the difference of arm means equals \(\beta+\varepsilon\mu_H\) when
      the contaminant \(H\) of \(F_1=(1-\varepsilon)\Phi+\varepsilon H\)
      has a mean \(\mu_H\): unbounded over the family, and undefined for
      a contaminant without a mean.
\item \emph{A robust decision without identification.} Let the declared
      family be
      \(\Theta_\Gamma=\{M(\beta,\gamma,F):\beta\in\mathbb R,\
      \gamma\in[0,\Gamma],\ F\in\mathcal F_\varepsilon\}\), with
      \(\Gamma\geq0\), \(0\leq\varepsilon<1\) and uncontaminated
      records, and suppose the clipped contrast has limit \(\alpha\).
      The members consistent with this limit have shift effects filling
      \([\alpha-\Gamma,\alpha]\). For the decision to adopt exactly when
      the shift effect exceeds \(c\), the conclusion that the shift
      effect exceeds \(c\), and with it the
      strict preference for adoption at every consistent member, is
      robust exactly when \(\alpha-\Gamma>c\), and the shift
      effect is identified within the consistent members only if
      \(\Gamma=0\). With separate record contaminations as in item 2 and
      limit \(\tau\), the consistent shift effects fill
      \([\tau-2c_*-\Gamma,\tau+2c_*]\), and adoption is robust exactly
      when \(\tau-2c_*-\Gamma>c\).
\item \emph{Robust is not correct.} For any \(\alpha\) and \(c\), the
      model \(M(\alpha-\gamma,\gamma,F)\) with \(\gamma>\alpha-c\) has
      the same observational law and a shift effect below \(c\). If
      \(\gamma>\Gamma\), it lies outside \(\Theta_\Gamma\), and the
      robust recommendation of item 4 is wrong for it.
\end{enumerate}
\end{hbtheorem}

\begin{hbexample}{def:p15-w02b-p15-070-separation-numbers}{The worked case with numbers}
Take \(F=\Phi\), \(\beta=0\), \(\gamma=0.4\) and \(k=1\). The two arms
are normal with centers \(0\) and \(0.4\), so the clipped contrast tends
to \(\alpha=0.4\) and remains exactly \(0.4\) under any common
contamination: a robust estimate of the observational contrast. The shift
effect is \(0\): setting \(X\) changes nothing, since every difference
between the arms comes from \(Z\).

Suppose an analyst adopts a treatment when its effect exceeds \(c=0.1\),
and declares that confounding can shift the contrast by at most
\(\Gamma=0.2\). Over that family the consistent effects are
\([0.2,0.4]\), and adoption is a robust decision although the effect is
not identified. The data-generating model above, with \(\gamma=0.4\),
lies outside the declared family; for it, adoption is wrong.
Robustness was relative to the declared range of confounding, and
nothing in the observational law tests that range (item 1).

If instead the observed limit \(\tau=0.4\) comes from records that may
have been contaminated separately in the two arms at
\(\varepsilon=0.05\), the bound \(c_*(0.05)<0.0772\) gives consistent
effects from below \(0.4-0.1542-0.2=0.0458\) up to above \(0.5542\).
Adoption is no longer robust: the lower end lies below \(0.1\). The
confounding margin \(\Gamma\) and the contamination margin \(2c_*\) add,
as in \cref{thm:p15-w02b-p15-070-margins}, because the two sources vary
separately and both attain their bounds.

In the randomized subclass \(\gamma=0\), the effect is identified, but a
record contamination of the treated arm with \(\varepsilon=0.05\) and a
contaminant \(H\) of mean \(20\) moves the difference of means from
\(\beta\) to \(\beta+1\), while the limit of the clipped contrast moves by
at most \(c_*(0.05)<0.0772\). Identification fixes the target; it does not choose a robust
way to estimate it.
\end{hbexample}

\begin{hbmethod}{def:p15-w02b-p15-070-which-claim}{Stating which robustness claim is made}
\begin{enumerate}
\item For a statistical claim, name the procedure, its target functional
      of the observational law, the neighborhood and the criterion
      (bias, limiting risk, coverage), and whether the bound is exact,
      asymptotic or approximate.
\item For a causal claim, name the model class and the intervention
      quantity, and state the identification result used. A procedure
      that is consistent across a class identifies its limit in that
      class (\cref{thm:p15-w02b-p15-070-consistency-identification});
      if the class does not identify the effect, report the set of
      effects consistent with the observational law under declared
      assumptions instead of a point value.
\item For a decision, name the actions, the declared family and the loss.
      Report whether a robustly optimal action exists. If it does not and
      a minimax or minimax-regret action is used, report it as a
      worst-case choice and state the loss normalization.
\item Combine margins only when each is valid throughout the declared
      family, and add them only after checking joint attainment.
\end{enumerate}
\end{hbmethod}

\section{Relationships, limitations and sources}
\label{sec:p15-w02b-p15-070-relations}

The threshold certificates of
\cref{thm:p15-approximation-certificates} and the reversal of
\cref{thm:p15-approximation-rational} are robustness statements over a
family indexed by an approximation error. The ranking and threshold
reversals of \cref{thm:p15-global-ranking} concern a family indexed by
input ranges, and \cref{thm:p15-w02b-p15-070-linear-ranking} extends its
exact criterion to every linear additive model and to a share threshold.
The causal boundary of \cref{thm:p15-global-causal-boundary} is the
sensitivity-index case of
\cref{thm:p15-w02b-p15-070-consistency-identification}: an index computed
from an observational law cannot identify an intervention response that
the law does not determine.
Huber notes, for a different pair of contaminated densities, that no
translation-invariant estimator can keep its bias below half the shift
separating them for both at once \cite[Sec.~7, p.~83]{Huber1964RobustLocation};
item 2 of \cref{thm:p15-w02b-p15-070-contamination-margin} is the
corresponding exact statement for the limit of the clipped estimate.
When an effect is not identified, Pearl describes bounding it by the range
of values compatible with the observed distribution, a range that more data
cannot narrow \cite[Sec.~3.5.3, p.~124]{Pearl2009CausalOverview}; the
interval \([\alpha-\Gamma,\alpha]\) of
\cref{thm:p15-w02b-p15-070-separation} is a bound of this kind, obtained
from a declared range for confounding rather than from the observational
law alone.

The results here are conditional on declared families. They do not say
which family is credible, which loss is appropriate, or whether an
asymptotic statement controls a finite sample. The contamination margins
use the known-scale clipped normal score; an unknown scale, another score
or another neighborhood requires its own displacement analysis
(\cref{sec:p15-huber-remaining-methods}).
